%! Author = chtompki
%! Date = 1/14/26
% Example LaTeX document for GP111 - note % sign indicates a comment
\documentclass[12pt]{amsart}
% Default margins are too wide all the way around. I reset them here
\setlength{\hoffset}{-.75in}
\setlength{\textwidth}{6.5in}
\setlength{\voffset}{-.5in}
\setlength{\textheight}{9.0in}
\usepackage{amsmath}
\usepackage{amssymb}
\usepackage{amsthm}
\usepackage{pgfplots}
\usepackage{enumerate}
\usepackage{tikz}
\usetikzlibrary{shadows,arrows,arrows.meta,shapes,positioning,calc}% Required for the Circle and Triangle arrow tips
\usepackage{setspace}
\usepackage{siunitx}

\theoremstyle{definition}
\newtheorem{definition}{Definition}

\theoremstyle{question}
\newtheorem{question}{Question}

\title{Seminar Summary:\\March 24, 2026 - Reed Ogrosky\\Some mathematical modeling projects\\ using differential equations\\ in fluid dynamics}
\author{Rob Tompkins}
\date{\today}


\begin{document}
    \maketitle
    \date{\today}
    \begin{onehalfspace}
        \section{Mathematical modeling of gas---liquid transport in human airways/lungs}
        The lecture began with a presentation of a model of the fluid mechanics of the movement of mucus around the lungs,
        which is a curious problem because there is quite a large free surface of liquid in the lungs due to the branches
        as well as the fluid itself being non-newtonian (having viscoelastic properties) in nature.
        We began by discussing the mechanics of the cilia that line the lungs and how they beat in rhythmic fashions
        to move the mucus up and out of the lungs generally.
        We were then posed with a quote:
        \begin{quote}
            ``It is awkward that in 2022 a comprehensive description of the organization and topography of mucus
            transport macroscopically under basal or stressed conditions is not available.'' --- Hill et al. (2022)
        \end{quote}


    \end{onehalfspace}
\end{document}
